% GENERATED FILE. Edit canonical lessons, then run npm run generate:books.
% canonical-source-hash: fnv1a64:beba0ccf04723623
% canonical-lessons: PA-C358-visha, PA-C358-ganit, PA-C358-vigian, PA-C358-itihas, PA-C358-parhauna

\chapter{What We Study}
\label{ch:pa-a2-what-we-study}
\hlchaptermodality{358}

\begin{chapteropening}
\textbf{By the end of this chapter:} \emph{I can name my school subjects and say who teaches them.}

Complete the last lesson of chapter 358: I can name my school subjects and say who teaches them.
\end{chapteropening}

\section[\texorpdfstring{\pa{ਵਿਸ਼ਾ} (vishā)}{ਵਿਸ਼ਾ (vishā)}]{\pa{ਵਿਸ਼ਾ} (vishā) — a school subject; a topic}

\label{lesson:PA-C358-visha}



\begin{quote}
Before the new one: say the Punjabi for a slate, then the Punjabi for a page.
\end{quote}

\subsection*{You'll want to know: \pa{ਵਿਸ਼ਾ}}
\textbf{\pa{ਵਿਸ਼ਾ}} — \emph{vishā} — \textquotedblleft{}a school subject; a topic\textquotedblright{}.

\begin{quote}
\pa{ਸਕੂਲ} \pa{ਵਿੱਚ} \pa{ਮੇਰੇ} \pa{ਛੇ} \pa{ਵਿਸ਼ੇ} \pa{ਹਨ।}
\end{quote}

\begin{quote}
\emph{sakūl vich mere chhe vishe han.}
\end{quote}

\begin{quote}
I have six subjects at school.
\end{quote}

A masculine noun.

\begin{cousinweb}
\textbf{\pa{ਵਿਸ਼ਾ}} (\emph{vishā}) is taken over from Sanskrit \emph{viṣaya}, \textquotedblleft{}sphere, scope\textquotedblright{}, a tatsama (Monier-Williams).
\end{cousinweb}

\begin{grammarlens}[title={What you've built}]
One more word for this chapter.
\end{grammarlens}

\subsection*{Your turn}
\begin{itemize}
\raggedright
  \item \emph{Say it:} \emph{vishā}
  \item \emph{Say it:} \emph{vishā}, once more
  \item \emph{Say it:} \emph{sakūl vich mere chhe vishe han.} — I have six subjects at school.
  \item \emph{Recall:} say the Punjabi for a slate, then the Punjabi for a page, then say \emph{vishā} again
\end{itemize}

\begin{tcolorbox}[breakable,colback=teal!4,colframe=teal!35!black,title={Before you move on}]
What does \pa{ਵਿਸ਼ਾ} mean? (\textquotedblleft{}A school subject; a topic\textquotedblright{}.) Say it once more.
\end{tcolorbox}
\section[\texorpdfstring{\pa{ਗਣਿਤ} (gaṇit)}{ਗਣਿਤ (gaṇit)}]{\pa{ਗਣਿਤ} (gaṇit) — mathematics}

\label{lesson:PA-C358-ganit}



\begin{quote}
Before the new one: say the Punjabi for a page, then the Punjabi for a school subject.
\end{quote}

\subsection*{You'll want to know: \pa{ਗਣਿਤ}}
\textbf{\pa{ਗਣਿਤ}} — \emph{gaṇit} — \textquotedblleft{}mathematics\textquotedblright{}.

\begin{quote}
\pa{ਗਣਿਤ} \pa{ਮੇਰੇ} \pa{ਲਈ} \pa{ਸੌਖਾ} \pa{ਹੈ।}
\end{quote}

\begin{quote}
\emph{gaṇit mere laī saukhā hai.}
\end{quote}

\begin{quote}
Mathematics is easy for me.
\end{quote}

A masculine noun.

\begin{cousinweb}
\textbf{\pa{ਗਣਿਤ}} (\emph{gaṇit}) is taken over from Sanskrit \emph{gaṇita}, \textquotedblleft{}reckoning, calculating\textquotedblright{}, a tatsama (Monier-Williams).
\end{cousinweb}

\begin{grammarlens}[title={What you've built}]
One more word for this chapter.
\end{grammarlens}

\subsection*{Your turn}
\begin{itemize}
\raggedright
  \item \emph{Say it:} \emph{gaṇit}
  \item \emph{Say it:} \emph{gaṇit}, once more
  \item \emph{Say it:} \emph{gaṇit mere laī saukhā hai.} — Mathematics is easy for me.
  \item \emph{Recall:} say the Punjabi for a page, then the Punjabi for a school subject, then say \emph{gaṇit} again
  \item \emph{Write it:} copy \textbf{\pa{ਗਣਿਤ}} once, checking the whole word
\end{itemize}

\begin{tcolorbox}[breakable,colback=teal!4,colframe=teal!35!black,title={Before you move on}]
What does \pa{ਗਣਿਤ} mean? (\textquotedblleft{}Mathematics\textquotedblright{}.) Say it once more.
\end{tcolorbox}
\section[\texorpdfstring{\pa{ਵਿਗਿਆਨ} (vigiān)}{ਵਿਗਿਆਨ (vigiān)}]{\pa{ਵਿਗਿਆਨ} (vigiān) — science}

\label{lesson:PA-C358-vigian}



\begin{quote}
Before the new one: say the Punjabi for a school subject, then the Punjabi for mathematics.
\end{quote}

\subsection*{You'll want to know: \pa{ਵਿਗਿਆਨ}}
\textbf{\pa{ਵਿਗਿਆਨ}} — \emph{vigiān} — \textquotedblleft{}science\textquotedblright{}.

\begin{quote}
\pa{ਮੇਰੀ} \pa{ਭੈਣ} \pa{ਵਿਗਿਆਨ} \pa{ਪੜ੍ਹਦੀ} \pa{ਹੈ।}
\end{quote}

\begin{quote}
\emph{merī bhaiṇ vigiān paṛhdī hai.}
\end{quote}

\begin{quote}
My sister studies science.
\end{quote}

A masculine noun.

\begin{cousinweb}
\textbf{\pa{ਵਿਗਿਆਨ}} (\emph{vigiān}) is taken over from Sanskrit \emph{vijñāna}, \textquotedblleft{}understanding, knowledge\textquotedblright{}, and also \textquotedblleft{}science\textquotedblright{}, a tatsama (Monier-Williams).
\end{cousinweb}

\begin{grammarlens}[title={What you've built}]
One more word for this chapter.
\end{grammarlens}

\subsection*{Your turn}
\begin{itemize}
\raggedright
  \item \emph{Say it:} \emph{vigiān}
  \item \emph{Say it:} \emph{vigiān}, once more
  \item \emph{Say it:} \emph{merī bhaiṇ vigiān paṛhdī hai.} — My sister studies science.
  \item \emph{Recall:} say the Punjabi for a school subject, then the Punjabi for mathematics, then say \emph{vigiān} again
  \item \emph{Write it:} copy \textbf{\pa{ਵਿਗਿਆਨ}} once, checking the whole word
\end{itemize}

\begin{tcolorbox}[breakable,colback=teal!4,colframe=teal!35!black,title={Before you move on}]
What does \pa{ਵਿਗਿਆਨ} mean? (\textquotedblleft{}Science\textquotedblright{}.) Say it once more.
\end{tcolorbox}
\section[\texorpdfstring{\pa{ਇਤਿਹਾਸ} (itihās)}{ਇਤਿਹਾਸ (itihās)}]{\pa{ਇਤਿਹਾਸ} (itihās) — history}

\label{lesson:PA-C358-itihas}



\begin{quote}
Before the new one: say the Punjabi for mathematics, then the Punjabi for science.
\end{quote}

\subsection*{You'll want to know: \pa{ਇਤਿਹਾਸ}}
\textbf{\pa{ਇਤਿਹਾਸ}} — \emph{itihās} — \textquotedblleft{}history\textquotedblright{}.

\begin{quote}
\pa{ਮੈਨੂੰ} \pa{ਇਤਿਹਾਸ} \pa{ਦੀ} \pa{ਹਰ} \pa{ਕਹਾਣੀ} \pa{ਪਸੰਦ} \pa{ਹੈ।}
\end{quote}

\begin{quote}
\emph{mainū̃ itihās dī har kahāṇī pasand hai.}
\end{quote}

\begin{quote}
I like every story from history.
\end{quote}

A masculine noun.

\begin{cousinweb}
\textbf{\pa{ਇਤਿਹਾਸ}} (\emph{itihās}) is taken over from Sanskrit \emph{itihāsa}, \textquotedblleft{}history, traditional accounts\textquotedblright{}, a tatsama (Monier-Williams).
\end{cousinweb}

\begin{grammarlens}[title={What you've built}]
One more word for this chapter.
\end{grammarlens}

\subsection*{Your turn}
\begin{itemize}
\raggedright
  \item \emph{Say it:} \emph{itihās}
  \item \emph{Say it:} \emph{itihās}, once more
  \item \emph{Say it:} \emph{mainū̃ itihās dī har kahāṇī pasand hai.} — I like every story from history.
  \item \emph{Recall:} say the Punjabi for mathematics, then the Punjabi for science, then say \emph{itihās} again
  \item \emph{Write it:} copy \textbf{\pa{ਇਤਿਹਾਸ}} once, checking the whole word
\end{itemize}

\begin{tcolorbox}[breakable,colback=teal!4,colframe=teal!35!black,title={Before you move on}]
What does \pa{ਇਤਿਹਾਸ} mean? (\textquotedblleft{}History\textquotedblright{}.) Say it once more.
\end{tcolorbox}
\section[\texorpdfstring{\pa{ਪੜ੍ਹਾਉਣਾ} (paṛhāuṇā)}{ਪੜ੍ਹਾਉਣਾ (paṛhāuṇā)}]{\pa{ਪੜ੍ਹਾਉਣਾ} (paṛhāuṇā) — to teach a subject; to give lessons}

\label{lesson:PA-C358-parhauna}



\begin{quote}
Before the new one: say the Punjabi for science, then the Punjabi for history.
\end{quote}

\subsection*{You'll want to know: \pa{ਪੜ੍ਹਾਉਣਾ}}
\textbf{\pa{ਪੜ੍ਹਾਉਣਾ}} — \emph{paṛhāuṇā} — \textquotedblleft{}to teach a subject; to give lessons\textquotedblright{}.

\begin{quote}
\pa{ਮੇਰੀ} \pa{ਮਾਤਾ} \pa{ਸਕੂਲ} \pa{ਵਿੱਚ} \pa{ਗਣਿਤ} \pa{ਪੜ੍ਹਾਉਂਦੀ} \pa{ਹੈ।}
\end{quote}

\begin{quote}
\emph{merī mātā sakūl vich gaṇit paṛhāundī hai.}
\end{quote}

\begin{quote}
My mother teaches mathematics at a school.
\end{quote}

Past (I …): \pa{ਮੈਂ} \pa{ਗਣਿਤ} \pa{ਪੜ੍ਹਾਇਆ}, \emph{maiṁ gaṇit paṛhāiā}.

\begin{cousinweb}
\textbf{\pa{ਪੜ੍ਹਾਉਣਾ}} (\emph{paṛhāuṇā}) is the causative partner of \emph{paṛhnā}, from Sanskrit \emph{paṭhati}, \textquotedblleft{}recites, repeats aloud\textquotedblright{}; CDIAL lists Punjabi \emph{paṛhāuṇā}, \textquotedblleft{}teaches\textquotedblright{}, under it (CDIAL 7712).
\end{cousinweb}

\begin{grammarlens}[title={What you've built}]
One more word for this chapter.
\end{grammarlens}

\subsection*{Your turn}
\begin{itemize}
\raggedright
  \item \emph{Say it:} \emph{paṛhāuṇā}
  \item \emph{Say it:} \emph{paṛhāuṇā}, once more
  \item \emph{Say it:} \emph{merī mātā sakūl vich gaṇit paṛhāundī hai.} — My mother teaches mathematics at a school.
  \item \emph{Recall:} say the Punjabi for science, then the Punjabi for history, then say \emph{paṛhāuṇā} again
  \item \emph{Write it:} copy \textbf{\pa{ਪੜ੍ਹਾਉਣਾ}} once, checking the whole word
\end{itemize}

\begin{tcolorbox}[breakable,colback=teal!4,colframe=teal!35!black,title={Before you move on}]
What does \pa{ਪੜ੍ਹਾਉਣਾ} mean? (\textquotedblleft{}To teach a subject; to give lessons\textquotedblright{}.) Say it once more.
\end{tcolorbox}
